\section{Evaluation Design}

This draft does not report experimental results. Instead, it specifies the evaluation that should be used once the simulator and scheduling algorithms are implemented.

\subsection{Scenarios}

The first experimental environment should be a single four-way intersection with a fixed conflict-zone partition. Vehicle streams can be generated synthetically before moving to real or simulator-derived trajectory data. The scenario factors should include:
\begin{itemize}
    \item traffic density: low, medium, and high arrival rates;
    \item vehicle mix: cars only; cars and buses; cars, buses, and trucks; and an emergency-vehicle scenario;
    \item priority setting: equal weights, moderate priority, strong priority, and hard emergency priority;
    \item prediction quality: perfect information, kinematic prediction, and noisy prediction.
\end{itemize}

\subsection{Metrics}

\begin{table}[t]
\centering
\caption{Planned Evaluation Metrics}
\label{tab:metrics}
\begin{tabular}{ll}
\toprule
Metric & Purpose \\
\midrule
Average delay & Overall efficiency \\
Weighted delay & Main priority-aware objective \\
Class-specific delay & Bus, truck, emergency service quality \\
Maximum delay & Starvation and fairness risk \\
Throughput & Vehicles served per unit time \\
Queue length & Congestion indicator \\
Stops/braking events & Comfort and energy proxy \\
Runtime per horizon & Real-time feasibility \\
Collision count & Must remain zero \\
Deadlock count & Must remain zero \\
\bottomrule
\end{tabular}
\end{table}

Class-specific reporting is essential. A method that lowers weighted delay by heavily delaying ordinary cars may be unacceptable unless the fairness threshold is satisfied. Similarly, a method that lowers average delay but ignores emergency vehicles should not be considered successful in priority-aware scenarios.

\subsection{Hypotheses}

The following hypotheses should be tested rather than assumed:
\begin{enumerate}
    \item Priority-aware scheduling reduces weighted delay and bus/emergency delay compared with equal-priority scheduling.
    \item Under low traffic density, FCFS or greedy policies may perform close to learning-based scheduling.
    \item Under high traffic density and mixed vehicle classes, PPO/GNN or hybrid scheduling may outperform FCFS and simple greedy methods.
    \item Fairness penalties reduce maximum delay for ordinary vehicles with an acceptable tradeoff in high-priority service.
    \item Receding-horizon scheduling with safety margins remains feasible under controlled prediction error.
\end{enumerate}

\subsection{Acceptance Criteria for the First Implementation}

The first implementation should be considered successful if it can:
\begin{itemize}
    \item generate feasible schedules with zero conflict-zone collisions and zero deadlocks in all tested scenarios;
    \item compute average delay, weighted delay, class-specific delay, fairness, throughput, and runtime;
    \item compare FCFS, greedy, and priority-greedy baselines before RL is added;
    \item validate small horizons against an exact optimization benchmark or clearly document why the benchmark is not yet available.
\end{itemize}

\placeholder{Insert experiment tables and plots after implementation. Do not add performance claims until the simulator and metrics are complete.}

